\documentclass[11pt]{article}
\usepackage[a4paper,margin=26mm]{geometry}
\usepackage{amsmath,amssymb,amsthm,booktabs,array,microtype}
\usepackage{xurl}
\usepackage[colorlinks=true,linkcolor=blue,urlcolor=blue,citecolor=blue]{hyperref}
\hypersetup{pdftitle={The exact seven-product closure degree is forty-seven},
  pdfauthor={Marcus Webb}}
\setlength{\emergencystretch}{2em}
\newtheorem{theorem}{Theorem}
\newtheorem{lemma}[theorem]{Lemma}
\newtheorem{proposition}[theorem]{Proposition}
\newcommand{\C}{\mathbb C}
\newcommand{\V}[1]{\C[x]_{\le #1}}
\newcommand{\Span}{\operatorname{span}}
\newcommand{\coeff}[1]{[x^{#1}]}
\newcommand{\norm}[1]{\left\lVert #1\right\rVert}
\title{The exact seven-product closure degree is forty-seven}
\author{Marcus Webb\\\small The University of Manchester}
\date{17 September 2026}

\begin{document}
\maketitle

\begin{abstract}
Starting from $1,x$ and allowing arbitrary complex linear combinations
for free, we determine the largest degree whose entire coefficient space
lies in the Zariski closure of outputs computable with seven products.
The answer is $47$. A $49$-parameter polynomial family has an exactly
certified contact of order $47$ with a nonsingular $48$-variable
coefficient differential. An inverse-function and scaling argument gives
coverage in the full degree-$128$ ambient space. A self-contained
dimension bound, together with irreducibility, proves that degree $48$
cannot be covered. The computational certificate uses rational arithmetic
and has been independently reconstructed.
\end{abstract}

\noindent\textbf{Assistance and review.}
This work was developed with substantial ChatGPT and Codex assistance.
Separate AI agents reviewed the mathematical argument and independently
implemented the certificate checks. Their reports and executable evidence
accompany this manuscript. These are informal AI-agent reviews, not
external human peer review or proof-kernel formalization.

\section{Statement and computational model}\label{sec:model}

Write $V_d=\V d$. Begin with the scalar polynomials $1,x\in\C[x]$.
Arbitrary complex linear combinations of available polynomials are free;
forming a product of two available polynomials costs one multiplication.
Let $\mathcal P_7\subseteq V_{128}$ be the possible outputs using at most
seven products, and define
\begin{equation}\label{eq:target}
 X_7=\overline{\mathcal P_7}^{\,Z}\subseteq V_{128},
 \qquad
 d_7=\max\{d\in\{0,\ldots,128\}:V_d\subseteq X_7\}.
\end{equation}
The closure is taken in all $129$ coefficients. In particular, an
approximating output may have higher degree than its limit. Discarding
high-degree coefficients is not a permitted computational operation.

\begin{theorem}\label{thm:main}
In the model \eqref{eq:target}, one has $d_7=47$.
\end{theorem}

MF-14 and Conjecture~13 of Jarlebring--Lorentzon~\cite{JL} asked whether
this maximum equals $42$. Colbrook~\cite{Colbrook} proved the degree-$42$
coverage statement with an exact coefficient-rank certificate. The
author's earlier degree-$44$ construction~\cite{Webb44} refuted the
equality to $42$ and left $44\le d_7\le47$. Theorem~\ref{thm:main}
determines the exact maximum. Its proof below is independent of those
lower-bound constructions and requires no auxiliary theorem for a
smaller multiplication budget.

The assertion concerns complex coefficient-space closure. It does not
assert exact representation of every degree-$47$ polynomial, real
coverage, or numerical stability of coefficient recovery. Substituting
a matrix for $x$ gives the corresponding matrix-polynomial evaluation
model, in which matrix products dominate the cost.

\section{A polynomial family in the seven-product closure}\label{sec:family}

Let $\alpha,b,c,\eta\in\C$ and set
\begin{equation}\label{eq:first-three}
 q_1=x^2,\qquad q_2=Q=x^4+\alpha x^3,\qquad
 q_3=R=Q^2+b x^2Q+c xQ+\eta x^3.
\end{equation}
For $k=4,5,6,7$, introduce the ordered list
\[
 L_k=(r_0,\ldots,r_{k-2})=(x,q_1,\ldots,q_{k-2})
\]
and define
\begin{equation}\label{eq:later-products}
 q_k=
 \left(q_{k-1}+\sum_{j=0}^{k-2}u_{k,j}r_j\right)
 \left(q_{k-1}+\sum_{j=0}^{k-2}v_{k,j}r_j\right).
\end{equation}
Finally, put
\begin{equation}\label{eq:output}
 \Phi(\theta)=e_0+e_1x+\sum_{k=1}^{7}e_{k+1}q_k.
\end{equation}
There are $4+(6+8+10+12)+9=49$ parameters. Their order is
\begin{equation}\label{eq:parameter-order}
 \theta=(\alpha,b,c,\eta,
          \mathbf u_4,\mathbf v_4,
          \mathbf u_5,\mathbf v_5,
          \mathbf u_6,\mathbf v_6,
          \mathbf u_7,\mathbf v_7,e_0,\ldots,e_8),
\end{equation}
with each vector ordered by increasing $j$. Table~\ref{tab:parameters}
specifies the zero-based indices used by the certificate. The map
$\Phi:\C^{49}\to V_{128}$ is polynomial, and all its coefficients
are integers as polynomials in $\theta$.

\begin{table}[htbp]
\centering
\begin{tabular}{@{}lll@{}}
\toprule
Parameter block & Entries in order & Zero-based indices\\
\midrule
Initial parameters & $\alpha,b,c,\eta$ & $0,1,2,3$\\
Fourth product & $\mathbf u_4;\ \mathbf v_4$ & $4$--$6;\ 7$--$9$\\
Fifth product & $\mathbf u_5;\ \mathbf v_5$ & $10$--$13;\ 14$--$17$\\
Sixth product & $\mathbf u_6;\ \mathbf v_6$ & $18$--$22;\ 23$--$27$\\
Seventh product & $\mathbf u_7;\ \mathbf v_7$ & $28$--$33;\ 34$--$39$\\
Output & $e_0,\ldots,e_8$ & $40$--$48$\\
\bottomrule
\end{tabular}
\caption{The complete parameter ordering of \eqref{eq:parameter-order}.}
\label{tab:parameters}
\end{table}

\begin{lemma}\label{lem:membership}
For every $\theta\in\C^{49}$, one has $\Phi(\theta)\in X_7$.
If $b\ne2$, then $\Phi(\theta)\in\mathcal P_7$.
\end{lemma}

\begin{proof}
The first two products are $x\,x$ and $x^2(x^2+\alpha x)$.
For $b\ne2$, set
\begin{equation}\label{eq:third-scalars}
 g=b-1,\qquad
 \sigma=\frac{\eta-c+\alpha g}{b-2},\qquad
 \rho=c-\sigma.
\end{equation}
The identity
\begin{equation}\label{eq:third-identity}
 R=(Q+\rho x+g x^2)(Q+\sigma x+x^2)-gQ-\rho\sigma x^2
\end{equation}
therefore computes $R$ with one additional product. Indeed, after the
two displayed subtractions and the algebraic substitution
$x^4=Q-\alpha x^3$, the right-hand side is
\[
 Q^2+(\rho+\sigma)xQ+(g+1)x^2Q
       +(\rho+g\sigma-g\alpha)x^3.
\]
The three relevant coefficients equal $c,b,\eta$, respectively.
Every factor and subtraction in \eqref{eq:third-identity} uses only
previously available polynomials. In particular, the expansion does not
charge the unavailable monomial $x^3$ as a free operation.

Each of \eqref{eq:later-products} costs one product, for a total of seven;
\eqref{eq:output} is free. Products and quotients of scalar parameters
in \eqref{eq:third-scalars} are allowed scalar coefficients. The degrees
of $q_1,\ldots,q_7$ are bounded by $2,4,8,16,32,64,128$.

The set $\{b\ne2\}$ is dense in $\C^{49}$, and the coefficient map
$\Phi$ is polynomial even at $b=2$. The inverse image $\Phi^{-1}(X_7)$
is closed and contains that dense open set; hence it is all of
$\C^{49}$. No assertion that $\Phi$ is dominant onto $X_7$ is needed.
\end{proof}

\section{An exact contact and rank certificate}\label{sec:certificate}

The frozen certificate file
\begin{center}
\path{source/experiments/lower_full49_seed.json}
\end{center}
specifies a center $\theta_c\in\mathbb Q(i)^{49}$, an ordered list
$I$ of $48$ distinct active indices, and a matrix
$B\in\mathbb Q(i)^{48\times48}$. The center denominator is $10^{70}$,
and the matrix denominator is $10^{60}$. The active list is reproduced
in Appendix~\ref{app:indices}; the sole fixed coordinate is
$\theta_{48}=e_8$.

Let $z\in\C^{48}$ denote the active coordinates, in the order $I$.
Keep the other coordinate fixed at its value in $\theta_c$, and write
$\iota(z)\in\C^{49}$ for the resulting full parameter vector. Define
\begin{equation}\label{eq:contact-equations}
 F(z)=\bigl(\coeff0\Phi(\iota(z)),\ldots,
             \coeff{47}\Phi(\iota(z))\bigr)^T
            -(0,\ldots,0,1)^T,
 \qquad J(z)=DF(z).
\end{equation}
Let $z_c$ be the active part of $\theta_c$.

For a complex number use $|w|_1=|\operatorname{Re}w|+
|\operatorname{Im}w|$. For vectors use the maximum of these norms,
and for matrices use the maximum row sum of entry norms. The latter
bounds the induced vector norm; also $|uv|_1\le |u|_1|v|_1$.

\begin{lemma}[Certified contact]\label{lem:contact}
There is a point $z_*$ with $\norm{z_*-z_c}\le10^{-45}$ such that
\[
 F(z_*)=0,\qquad \det J(z_*)\ne0.
\]
The associated parameter $b$ is different from $2$ throughout this ball.
Thus the untruncated polynomial $\Phi(\iota(z_*))$ has coefficients
$0,\ldots,0,1$ in degrees $0,\ldots,47$, and its selected coefficient
differential has rank $48$.
\end{lemma}

\begin{proof}
The accompanying exact verifier evaluates the specified Gaussian-rational
data and proves the following rational inequalities, with
$r=10^{-45}$:
\begin{align}
 \varepsilon:=\norm{BF(z_c)}&<10^{-69},&
 \delta:=\norm{I-BJ(z_c)}&<10^{-57},\label{eq:certificate-residual}\\
 \norm B&<4000,& H&<2\cdot10^{12},\label{eq:certificate-hessian}\\
 q:=\delta+\norm B\,Hr&<10^{-29},&
 \varepsilon+qr&<r.\label{eq:certificate-contraction}
\end{align}
It also checks
\begin{equation}\label{eq:b-separation}
 |\operatorname{Re}(b_c)-2|>r.
\end{equation}
Here $H$ is a uniform aggregate second-derivative bound on the closed
radius-$r$ ball, as justified below. All these inequalities are evaluated
using integers and rational numbers. Printed floating-point values are
diagnostics, not certificate inputs.

For clarity, define the coefficient norm of a polynomial
$a(x)=\sum_j a_jx^j$ by $\norm a_{\mathrm{coef}}=\sum_j|a_j|_1$.
At every node of the full polynomial circuit propagate a triple
$(v,d,h)$ bounding its coefficient norm, the sum of the norms of its
$48$ first parameter derivatives, and the sum of the norms of all its
$48^2$ \emph{ordered} second parameter derivatives. Addition adds these
triples. Multiplication uses
\begin{equation}\label{eq:majorant-product}
 (v,d,h)\cdot(w,e,k)
       =(vw,\ dw+ve,\ hw+2de+vk).
\end{equation}
An active scalar parameter has initial triple
$(|\theta_{c,j}|_1+r,1,0)$; a fixed scalar parameter has
$(|\theta_{c,j}|_1,0,0)$; a fixed monomial has $(1,0,0)$.
The product rule and submultiplicativity of
$\norm{\cdot}_{\mathrm{coef}}$ prove these bounds inductively.
The final value of $h$ is the bound $H$ in
\eqref{eq:certificate-hessian}. In particular, it bounds the aggregate
second derivatives of every output coefficient. Integration along a
line segment in the convex ball gives
\begin{equation}\label{eq:jacobian-lipschitz}
 \norm{J(z)-J(z_c)}\le Hr.
\end{equation}

Consider $G(z)=z-BF(z)$. Equations
\eqref{eq:certificate-contraction}--\eqref{eq:jacobian-lipschitz} give
\[
 \sup_{\norm{z-z_c}\le r}\norm{DG(z)}\le q<\tfrac12,
 \qquad
 \norm{G(z)-z_c}\le\varepsilon+qr<r.
\]
The closed ball is complete, so Banach's fixed point theorem supplies
a fixed point $z_*$. Since $\delta<1$, the matrix $BJ(z_c)$, and hence
$B$, is invertible; thus $BF(z_*)=0$ implies $F(z_*)=0$. Moreover,
$\norm{I-BJ(z_*)}\le q<1$, so the Neumann series proves that
$BJ(z_*)$, and therefore $J(z_*)$, is invertible. Finally,
\eqref{eq:b-separation} keeps $b$ different from $2$ throughout the ball.
\end{proof}

\subsection{Exact arithmetic and independent reconstruction}

The author-side check is implemented in
\path{source/experiments/lower_full49_certify.py}, with Gaussian-integer
helpers in \path{source/experiments/lower_contact_certify.py}.
Value and Jacobian calculations may use $\C[x]/(x^{48})$, since this
computes exactly the equations \eqref{eq:contact-equations} and their
derivatives. The argument itself uses the untruncated $\Phi$ in
$V_{128}$; no truncated polynomial is asserted to be a circuit output.

The parameter-degree bounds at $q_1,Q,R,q_4,q_5,q_6,q_7,\Phi$ are
$0,1,2,4,8,16,32,33$. Consequently the value and derivative numerators
can share denominator $10^{2310}$. Derivative injection rescales
numerators for differentiation with respect to the actual rational
parameters, rather than their integer encodings. Matrix products, norms
and asserted inequalities use Gaussian integers and exact rational
arithmetic. The Hessian majorant traverses the full circuit and does not
discard coefficients.

The independent implementation
\begin{center}
\path{source/reviews/full49-degree47-fresh-check.py}
\end{center}
imports no author implementation. It evaluates all $129$ coefficients and constructs all
$49$ derivative columns by a reverse pass through the product gates.
It reproduces the exact preconditioned residual, inverse defect and
inverse norm. Its separate positive-polynomial majorant, using the
verified bound $4$ for every parameter norm, gives
\begin{equation}\label{eq:independent-hessian}
 H_{\mathrm{ind}}=14225297189516125194807607250462624.
\end{equation}
Although coarser, this bound still proves $q_{\mathrm{ind}}<6\cdot10^{-8}$
and $\varepsilon+q_{\mathrm{ind}}r<r$ at the same radius. That checker
also independently verifies $b\ne2$. Thus the existence and rank claims
do not depend on the author's derivative or Hessian implementation.
The numerical search and Decimal refinement are discovery records only;
the frozen rational seed suffices for reproduction.

\section{Full coefficient-space coverage}\label{sec:coverage}

\begin{proposition}\label{prop:lower}
Every polynomial in $V_{47}$ belongs to $X_7$.
\end{proposition}

\begin{proof}
Fix any monic polynomial $p(x)=\sum_{j=0}^{47}p_jx^j$ of degree $47$.
By Lemma~\ref{lem:contact}, the selected coefficient map is locally
invertible at $z_*$. The complex inverse function theorem therefore
gives a holomorphic curve $z(t)\to z_*$, with the remaining parameter
fixed, such that
\begin{equation}\label{eq:scaled-target}
 \coeff j\Phi(\iota(z(t)))=p_jt^{47-j},\qquad j=0,\ldots,47.
\end{equation}
At $t=0$ the prescribed coefficient vector is exactly
$(0,\ldots,0,1)$, so the local inverse applies for all sufficiently
small $t$. In particular $z(t)$, and hence all coefficients of
$\Phi(\iota(z(t)))$, remain bounded near zero.

For $t\ne0$ set
\begin{equation}\label{eq:full-limit}
 f_t(x)=t^{-47}\Phi(\iota(z(t)))(tx).
\end{equation}
Input substitution $x\mapsto tx$ and scalar output multiplication
preserve the computational model and its closure. Thus $f_t\in X_7$.
In fact, after shrinking the parameter neighbourhood if necessary,
Lemma~\ref{lem:membership} and $b\ne2$ give $f_t\in\mathcal P_7$.
By \eqref{eq:scaled-target}, its coefficients through degree $47$
are exactly those of $p$. For every $j=48,\ldots,128$, its coefficient
is $t^{j-47}$ times a bounded holomorphic function, and therefore tends
to zero. There are no coefficients above degree $128$.

Consequently $f_t\to p$ in the full space $V_{128}$. An algebraic set
over $\C$ is Euclidean closed, so $p\in X_7$. Scalar output multiplication
gives every nonzero leading coefficient. Finally, any lower-degree
polynomial is a limit of degree-$47$ polynomials, for example of
$p+\epsilon x^{47}$, and closedness includes it as well.
\end{proof}

This proof uses contact and scaling in addition to coefficient rank.
Dominance of a coefficient projection alone would not establish the
full-space closure assertion.

\section{A self-contained matching upper bound}\label{sec:upper}

\begin{lemma}\label{lem:dimension}
The variety $X_7$ is irreducible and has dimension at most $49$.
\end{lemma}

\begin{proof}
At each gate of a universal circuit, multiply two arbitrary linear
combinations of all previously available polynomials, and take a final
arbitrary linear combination. These coefficients form one affine
parameter space, and the output map into $V_{128}$ is polynomial.
Circuits using fewer than seven products can be padded by zero products.
Therefore its image closure is exactly $X_7$, which is irreducible.

It suffices to bound the output dimension on the nonempty open subset
where the successive product degrees are $2,4,\ldots,128$. This subset
is dense in the universal affine parameter space, and its image has
the same closure. Let $W_k$ be the space of available polynomials after
$k$ products; generically $\dim W_k=k+2$.

After the first product, a free change of basis gives $W_1=V_2$.
After the second, subtract lower-degree terms and normalize the leading
coefficient to obtain
\[
 W_2=\Span\{1,x,x^2,Q\},\qquad Q=x^4+\alpha x^3.
\]
This family has dimension at most one. Denote by $W_2W_2$ the linear
span of its pairwise products. Then
\begin{equation}\label{eq:product-span}
 W_2W_2=V_6+\C Q^2,\qquad \dim(W_2W_2)=8.
\end{equation}
Indeed, products of $V_2$ span $V_4$; then $xQ$ and $x^2Q$ add
degrees $5$ and $6$, while $Q^2$ adds degree $8$. These span every
pairwise product. Since $1\in W_2$, the space $W_2$ is contained in
$W_2W_2$. The third product determines a line in the four-dimensional
quotient $(W_2W_2)/W_2$. Hence the family of possible $W_3$ spaces
has dimension at most $1+3=4$.

For a fixed $n$-dimensional available space $W$ containing $1$, the
next factors $(u,v)\in W^2$ determine $\Span(W,uv)$. On the generic
set where $uv\notin W$ and both factors are nonconstant, each fiber
of this space-valued map has dimension at least four. Specifically,
\begin{equation}\label{eq:factor-gauges}
 (u,v)\longmapsto(au+b,cv+d),\qquad a,c\in\C\setminus\{0\},
 \quad b,d\in\C,
\end{equation}
preserves the enlarged space, because the transformed product is
$ac\,uv+ad\,u+bc\,v+bd$. Distinct quadruples in
\eqref{eq:factor-gauges} give distinct pairs when $u,v$ are nonconstant.
The fiber-dimension bound thus shows that one product increases the
dimension of a family of available spaces by at most $2n-4$.

For the four remaining products, $n=5,6,7,8$, so the increments are
$6,8,10,12$. The family of $W_7$ spaces therefore has dimension at
most $4+6+8+10+12=40$. An arbitrary output in the nine-dimensional
space $W_7$ adds at most nine dimensions. This gives
$\dim X_7\le40+9=49$.

The argument bounds a dense generic image and then its closure. It
does not identify limiting multiplication spaces with products of
limiting spaces, or rule out boundary circuits using degree profiles.
\end{proof}

\begin{proof}[Proof of Theorem~\ref{thm:main}]
Proposition~\ref{prop:lower} gives $d_7\ge47$. If $V_{48}\subseteq X_7$,
then $\dim V_{48}=49$ and Lemma~\ref{lem:dimension} force
$\dim X_7=49$. A proper closed subset of an irreducible variety has
smaller dimension, so this containment would imply $X_7=V_{48}$.
But seven successive squarings compute $x^{128}$, which belongs to
$\mathcal P_7$ and is outside $V_{48}$. Thus $V_{48}\nsubseteq X_7$,
and every $V_d$ with $d\ge48$ is excluded. Hence $d_7=47$.
\end{proof}

\section{Reproduction and evidence scope}\label{sec:reproduction}

From the accompanying contribution directory
\begin{center}
\path{references/webb-mf14-degree47-2026-09-17}
\end{center}
run
\begin{verbatim}
python3 verification/verify.py
\end{verbatim}
The contribution README records the frozen files, hashes, exact outputs
and separate review reports. The rational seed, rather than the numerical
search trajectory, is the input needed to reproduce the existence and
rank certificate. The two implementations check the finite arithmetic
claims; the analytic and algebraic arguments in this manuscript have
separate informal reviews. No Lean verification, proof-kernel checking,
or external human peer review is claimed.

\appendix
\section{Active-coordinate order}\label{app:indices}

Relative to the zero-based ordering in Table~\ref{tab:parameters}, the
certificate's active coordinates are listed below in the order used for
Jacobian columns and the rows of its preconditioner:
\[
\begin{array}{rrrrrrrrrrrr}
4&3&14&18&19&5&34&36&35&10&25&0\\
15&20&23&22&28&29&8&45&37&40&41&43\\
42&44&7&46&30&11&1&24&21&2&31&16\\
9&47&6&27&12&38&39&26&17&33&13&32.
\end{array}
\]
All indices except $48$ occur exactly once. The omitted coordinate
$e_8$ is held at the exact Gaussian-rational value in the seed file.
The preconditioner's columns correspond to coefficient equations in
increasing degree order $0,\ldots,47$.

\begin{thebibliography}{9}
\bibitem{JL}
E. Jarlebring and G. Lorentzon,
\emph{The polynomial set associated with a fixed number of matrix-matrix
multiplications}, arXiv:2504.01500v3, Theorem~10 and Conjecture~13.
\url{https://arxiv.org/html/2504.01500v3}.

\bibitem{Colbrook}
M. J. Colbrook,
\emph{An exact degree-42 coverage certificate for seven polynomial
multiplications}, 11 September 2026, Theorem~1.
Archived in the OpenProblemsInNLA reference collection,
\href{../colbrook-matrix-functions-2026-09-11/manuscripts/MF-14.pdf}
{\texttt{colbrook-matrix-functions-2026-09-11}}.

\bibitem{Webb44}
M. Webb,
\emph{Degree forty-four in the closure of seven-product polynomial
evaluations: A negative resolution of MF-14}, 17 September 2026,
Theorem~1 and Corollary~3.
Archived in the OpenProblemsInNLA reference collection,
\href{../webb-mf14-degree44-2026-09-17/proof.pdf}
{\texttt{webb-mf14-degree44-2026-09-17}}.
\end{thebibliography}
\end{document}
